\clearpage
\VerdierPhysicalPageAnchor{0232}%
\begingroup\small\setlength{\parskip}{0pt}\setlength{\emergencystretch}{1.5em}
\setlength{\abovedisplayskip}{4pt plus 2pt minus 2pt}%
\setlength{\belowdisplayskip}{4pt plus 2pt minus 2pt}%

\VerdierSemanticAnchor{ch4d-4-1-7-double-complex-case}
Supposons donc que $X$ soit un double complexe et que $m=1$. Pour
démontrer a), il suffit de montrer que, pour tout $(p,q,r)$ de
$\mathcal{F}\ell_2(\widetilde{\mathbb Z})$, le diagramme de
$\operatorname{comp}(\mathcal A)$ :

\VerdierSemanticAnchor{ch4d-4-1-7-6}
\begin{equation}
\resizebox{0.98\linewidth}{!}{$
\displaystyle
\int_2{}^{1}\!X(p,q)\longrightarrow
\int_2{}^{1}\!X(p,r)\longrightarrow
\int_2{}^{1}\!X(q,r)
\xrightarrow{{}^{1}\!d(p,q,r)}
\int_2T_1\bigl({}^{1}\!X(p,q)\bigr)
=\left(\int_2{}^{1}\!X(p,q)\right)[1]
$}
\tag{4.1.7.6}
\end{equation}

définit dans $K(\mathcal A)$ un triangle distingué (remarquons que
l'isomorphisme de commutation aux translations
$\int_2T_1({}^{1}\!X(p,q))\xrightarrow{\sim}
(\int_2{}^{1}\!X(p,q))[1]$ est dans ce cas l'identité). Les autres
assertions contenues dans a) sont évidentes. Pour prouver que le diagramme
(4.1.7.6) définit dans $K(\mathcal A)$ un triangle distingué, il suffit,
d'après la propriété (TRII) (chap.~II, 1.1.1), de montrer que le diagramme :

\VerdierSemanticAnchor{ch4d-4-1-7-7}
\begin{equation}
\resizebox{0.96\linewidth}{!}{$
\displaystyle
\left(\int_2{}^{1}\!X(q,r)\right)[-1]
\xrightarrow{-{}^{1}\!d(p,q,r)[-1]}
\int_2{}^{1}\!X(p,q)\longrightarrow
\int_2{}^{1}\!X(p,r)\longrightarrow
\int_2{}^{1}\!X(q,r)
$}
\tag{4.1.7.7}
\end{equation}

définit dans $K(\mathcal A)$ un triangle distingué. On constate alors
immédiatement que le complexe simple $\int_2{}^{1}\!X(p,r)$ est égal au
cône du morphisme $-{}^{1}\!d(p,q,r)[-1]$ (chap.~I, 3.1.2) et que les
morphismes
$\int_2{}^{1}\!X(p,q)\longrightarrow\int_2{}^{1}\!X(p,r)$,
$\int_2{}^{1}\!X(p,r)\longrightarrow\int_2{}^{1}\!X(q,r)$, figurant dans
(4.1.7.7), sont les morphismes introduits dans l'étude du foncteur cône
(chap.~I, 3.2.2.2 et 3.2.2.4). Par suite, le triangle de $K(\mathcal A)$
défini par le diagramme (4.1.7.7) est par définition (chap.~I, 3.3.1) un
triangle distingué. Démontrons maintenant b). En utilisant les notations
introduites en (4.1.6.5), on vient de montrer que dans le diagramme :

\VerdierSemanticAnchor{ch4d-4-1-7-8}
\begin{equation}
{}^{1}\!Y(q,r)[-1]
\xrightarrow{-{}^{1}\!d(-\infty,q,r)}
{}^{1}\!Y(q)\longrightarrow{}^{1}\!Y(r)\longrightarrow{}^{1}\!Y(q,r),
\tag{4.1.7.8}
\end{equation}

le complexe ${}^{1}\!Y(r)$ apparaît comme le cône du morphisme
$-{}^{1}\!d(-\infty,q,r)$ (poser $p=-\infty$ dans le diagramme
(4.1.7.7)). De plus (4.1.1), l'objet $\mathcal Y(q,r)$ a été défini comme
étant le cône du morphisme :

\VerdierSemanticAnchor{ch4d-4-1-7-9}
\begin{equation}
{}^{1}\!Y(q)\longrightarrow{}^{1}\!Y(r).
\tag{4.1.7.9}
\end{equation}

On a alors défini, dans cette situation (chap.~I, 3.2.5.2), un homotopisme
(définissant donc un isomorphisme dans $K(\mathcal A)$) :

\VerdierSemanticAnchor{ch4d-4-1-7-10}
\begin{equation}
J(q,r)\colon\mathcal Y(q,r)\longrightarrow{}^{1}\!Y(q,r).
\tag{4.1.7.10}
\end{equation}

\VerdierSemanticAnchor{ch4d-4-1-7-proof-end}
On laisse au lecteur le soin de vérifier que cet homotopisme est fonctoriel
par rapport aux indices $(q,r)$ de
$\mathcal{F}\ell_1(\widetilde{\mathbb Z})$ et par rapport au double
complexe $X$ et qu'il définit un isomorphisme canonique de l'objet spectral
$\bigl(\mathcal Y(p,q),\delta(p,q,r)\bigr)$ sur l'objet spectral
$\bigl({}^{1}\!Y(p,q),{}^{1}\!d(p,q,r)\bigr)$.
\endgroup
